\section{Verification and reuse}\label{sec:verification}

The authoritative proofs of the results marked \textsf{Lean} are in
the memo's \texttt{lean/CircularMemo/} library. The English outlines
describe those proofs for convenience. The general partition identity,
the general acyclic-orientation argument, unequal alphabets, and the
sharpness examples have written proofs; they are not claimed to be
formalized by the three-part inequality alone.

The library requires the fourth paper by the relative path
\texttt{../../lean}; that paper in turn requires the earlier shared
libraries. Clone the whole repository, preserving sibling directories.
Lean and Mathlib remain pinned to v4.34.1. No new axiom, proof placeholder,
or kernel-bypassing decision procedure is used. The permitted axiom set
is \texttt{propext}, \texttt{Classical.choice}, and \texttt{Quot.sound}.

From the memo directory, run \texttt{python3 tools/lean\_index.py}.
It generates this declaration index and \texttt{lean/PaperAxioms.lean}
from the same formal-status annotations. In \texttt{lean/}, run
\texttt{lake build} and \texttt{lake env lean PaperAxioms.lean}.
On a fresh clone first fetch the pinned dependency cache with
\texttt{lake exe cache get}. In \texttt{paper/}, run
\texttt{latexmk -pdf main.tex}.

The main paper's original declarations and certificate remain intact.
The memo's old-range corollary derives that range directly from the
stronger theorem, including the empty case.
Removing the old certificate from its original module is deferred until
a separate dependency and publication review.
